\documentclass[11pt]{article}
\usepackage[utf8]{inputenc}
\usepackage[T1]{fontenc}
\usepackage{geometry}
\geometry{margin=1in}
\usepackage{lmodern}
\usepackage{amsmath,amssymb,amsthm,mathtools}
\usepackage{bm}
\usepackage{physics}
\usepackage{graphicx}
\usepackage{caption}
\usepackage{subcaption}
\usepackage{booktabs}
\usepackage{microtype}
\usepackage{enumitem}
\usepackage{hyperref}
\hypersetup{colorlinks=true, linkcolor=blue, citecolor=blue, urlcolor=blue}
\graphicspath{{./}}

\title{\textbf{Temporal Foam and Envelopes (TFE):\\
Emergent Space, Measurement, and Superdeterministic Correlations}\\[4pt]
\large A Comprehensive Proposal with Simulations}
\author{Anonymous (you can add your name here)}
\date{\today}

\newtheorem{definition}{Definition}
\newtheorem{proposition}{Proposition}
\newtheorem{theorem}{Theorem}

\newcommand{\HH}{\mathcal{H}}
\newcommand{\VV}{\mathcal{V}}
\newcommand{\LL}{\mathcal{L}}
\newcommand{\Keff}{K_{\mathrm{eff}}}
\newcommand{\EE}{\mathbb{E}}
\newcommand{\PP}{\mathbb{P}}

\begin{document}
\maketitle

\begin{abstract}
We develop a self-contained framework in which ``space'' is not fundamental but emerges from a fully connected temporal substrate (\emph{temporal foam}). Structure is imposed by \emph{envelopes}, represented by operators that dephase or confine dynamics, yielding effective locality, finite signalling speed, and one-dimensional preferred paths we call \emph{timelines}. Particles are excitations that propagate at constant intrinsic speed along timelines; all apparent speeding, slowing, and reversal in a spatial readout are geometric projection effects. We show how measurement corresponds to envelope-induced phase locking (einselection), how mass arises from spectral gaps, transverse confinement, or internal mixing, and how a superdeterministic variant arises naturally from clock-phase synchronization in the foam, violating measurement independence while preserving locality and no-signalling. We present simple simulations: motion on line/circle/spiral timelines; a branch node with a unitary splitter and measurement-induced locking; and a CHSH test under the clock-phase model exhibiting $S(\varepsilon)=(1-\varepsilon)2\sqrt2+2\varepsilon$. We conclude with testable predictions and a roadmap for further work.
\end{abstract}

\tableofcontents

\section{Introduction}

Nonlocal quantum correlations, decoherence, and the success of holographic ideas suggest that spacetime locality may be emergent from deeper informational structure. We propose a compact formalism---\emph{Temporal Foam + Envelopes (TFE)}---that makes this idea concrete and readily simulable. The central ingredients are:
\begin{enumerate}[label=(\roman*)]
\item a fully connected ``foam'' of micro-events with complex couplings;
\item \emph{envelopes} (operators) that sculpt dynamics by dephasing or confining;
\item emergent coordinates extracted spectrally from the envelope-filtered couplings;
\item timelines (1-D subspaces) supporting constant intrinsic-speed propagation;
\item measurement as strong envelope-induced phase locking (einselection);
\item a superdeterministic variant realized as clock-phase synchronization in the foam.
\end{enumerate}

\section{The TFE Formalism}
\subsection{Primitives}
Let $\VV=\{1,\dots,N\}$ be micro-events and $\HH=\mathrm{span}\{|i,\sigma\rangle: i\in\VV,\ \sigma\in\Sigma\}$ the Hilbert space (internal labels $\sigma$ are optional). A complex fully connected kernel $K\in\mathbb{C}^{N\times N}$ is split into its Hermitian and skew parts
\begin{equation}
H=\tfrac12(K+K^\dagger),\qquad A=\tfrac{1}{2i}(K-K^\dagger).
\end{equation}
The intrinsic time parameter is $\tau\in\mathbb{R}$.

\subsection{Open-system dynamics with envelopes}
States evolve under a GKLS generator
\begin{equation}\label{eq:Lindblad}
\partial_\tau \rho=-\frac{i}{\hbar}[H,\rho]+\sum_\alpha \gamma_\alpha\Big(L_\alpha\rho L_\alpha^\dagger-\tfrac12\{L_\alpha^\dagger L_\alpha,\rho\}\Big),
\end{equation}
where the \emph{envelope operators} $L_\alpha$ encode structure. Two canonical types:
\begin{itemize}
\item \textbf{Dephasing envelopes} $L_\alpha=E_\alpha=E_\alpha^\dagger\ge0$ (phase locking, einselection).
\item \textbf{Confining envelopes} $L_\alpha=\sqrt{V_\alpha}$ (resistance to leave preferred sets).
\end{itemize}

\subsection{Timelines and resistance to leave them}
A \emph{timeline} is a smooth family of rank-1 projectors $\{P(s)\}$, with $s$ arclength. Leakage is penalized by
\begin{equation}\label{eq:timeline}
L_{\mathcal T}(s)=\sqrt{\kappa}\,\big(I-P(s)\big),\qquad \kappa>0.
\end{equation}
Inside $\mathrm{im}\,P(s)$, propagation is at constant intrinsic speed $ds/d\tau=c$. Apparent nonuniformities in a spatial \emph{readout} $X=\langle \hat x,\gamma(s)\rangle$ are purely geometric.

\subsection{Branch nodes and unitary splitters}
At a node $s_\*$ where a timeline splits into channels $a=1,\dots,m$ with projectors $P_a$, the outgoing amplitudes are related by a unitary $S\in U(m)$:
\begin{equation}
\Psi_\text{out}=S\,\Psi_\text{in}.
\end{equation}
Different leakage penalties $\kappa_a$ implement tunable splitting ratios.

\subsection{Emergent locality and speed limit}
Define an envelope-filtered effective adjacency $\Keff$ (symmetric, nonnegative), and the Laplacian $L_\text{diff}=D-\Keff$. The leading eigenvectors of $L_\text{diff}$ provide an embedding $\Phi:i\mapsto x(i)\in\mathbb{R}^d$ which serves as emergent coordinates. When $\Keff$ is banded in $x$, the unitary part of~\eqref{eq:Lindblad} reduces to a local PDE,
\begin{equation}\label{eq:localPDE}
i\hbar\,\partial_\tau \psi(x)\approx -\frac{\hbar^2}{2\mu}\,\Delta_g \psi(x)+U(x)\psi(x),
\end{equation}
and Lieb--Robinson-type bounds yield a finite maximal group velocity, identified with $c$.

\section{Photon Sector and Timeline Geometry}
Choose a massless band with linear dispersion, $E(k)\approx \hbar c \|k\|$, i.e. $H_\text{ph}\approx \hbar c\,\sqrt{-\Delta_g}$. Confinement to a 1-D curve $x=x_{\mathcal T}(s)$ is enforced by a strong transverse envelope.

\paragraph{Readout kinematics.} For a readout $X=\Re$, unit-speed motion along the curve yields
\begin{equation}
\dot X(\tau)=c\,\frac{dX}{ds} = c\,\langle T(s),\hat x\rangle,\qquad T=\partial_s x_{\mathcal T}.
\end{equation}
Thus $X(t)$ speeds up, slows, and reverses solely from the tangent's projection.

\subsection{Line, circle, and spiral (simulated)}
We simulate three shapes (figs.~\ref{fig:paths}--\ref{fig:spiral}):
\begin{itemize}
\item \textbf{Line:} $x_{\mathcal T}(s)=(s,0)$, $X=s$ and $\dot X=c$.
\item \textbf{Circle:} $x_{\mathcal T}(s)=(R\cos\frac{s}{R},R\sin\frac{s}{R})$, $X=R\cos\frac{s}{R}$ and $\dot X=-c\sin\frac{s}{R}$ (reverses at turning points).
\item \textbf{Archimedean spiral:} $r=a+b\theta$, $X=r\cos\theta$; reversals persist with growing amplitude and net drift each $2\pi$.
\end{itemize}

\begin{figure}[h!]
\centering
\includegraphics[width=.9\textwidth]{line_path.png}\\[3pt]
\includegraphics[width=.49\textwidth]{line_X_t.png}\hfill
\includegraphics[width=.49\textwidth]{line_vX_t.png}
\caption{Linear timeline: path (top), readout $X(t)$, and apparent speed $\dot X(t)$.}
\label{fig:paths}
\end{figure}

\begin{figure}[h!]
\centering
\includegraphics[width=.9\textwidth]{circle_path.png}\\[3pt]
\includegraphics[width=.49\textwidth]{circle_X_t.png}\hfill
\includegraphics[width=.49\textwidth]{circle_vX_t.png}
\caption{Circular timeline: periodic reversals in the readout arise from projection geometry.}
\label{fig:circle}
\end{figure}

\begin{figure}[h!]
\centering
\includegraphics[width=.9\textwidth]{spiral_path.png}\\[3pt]
\includegraphics[width=.49\textwidth]{spiral_X_t.png}\hfill
\includegraphics[width=.49\textwidth]{spiral_vX_t.png}
\caption{Spiral timeline: reversals with growing amplitude and net drift.}
\label{fig:spiral}
\end{figure}

\section{Mass in TFE}
We present three equivalent definitions.
\paragraph{(i) Spectral gap (primary):}
\begin{equation}
m c^2 := E(0)=\hbar \omega(0),\qquad
E(k)\approx \sqrt{(mc^2)^2+(\hbar c\,\|k\|)^2}.
\end{equation}
\paragraph{(ii) Confinement/waveguide:} For transverse eigenvalues $\lambda_n$ of the tube,
\begin{equation}
E_n(k)=\sqrt{(\hbar c)^2\|k\|^2+(\hbar c)^2\lambda_n}\quad\Rightarrow\quad
m_n c^2=\hbar c\,\sqrt{\lambda_n}.
\end{equation}
\paragraph{(iii) Internal mixing (Dirac-like):} If an envelope mixes internal labels with strength $M$,
\begin{equation}
H \approx \hbar c\, \bm{\alpha}\!\cdot\!\bm{k}+ \beta\, M,\qquad m=M/c^2.
\end{equation}

\section{Measurement as Envelope-Induced Phase Locking}
Strong dephasing envelopes $L_\text{obs}$ select pointer states and suppress coherences. In a two-time (pre/post) description, probabilities follow the Aharonov--Bergmann--Lebowitz rule while einselection fixes the pointer basis.

\section{Branch Nodes and Measurement Locking (Simulation)}
We inject a Gaussian wavepacket into a node with a 50/50 unitary splitter $S$ feeding (A) a line and (B) a circle. Without measurement, ensemble $X(t)$ is a weighted sum of the two geometries. Turning on strong which-path dephasing at $t_\text{lock}$ locks the readout to a single pointer path (fig.~\ref{fig:branch}).
\begin{figure}[h!]
\centering
\includegraphics[width=.9\textwidth]{branch_geometry.png}\\[3pt]
\includegraphics[width=.49\textwidth]{X_ensemble_no_measurement.png}\hfill
\includegraphics[width=.49\textwidth]{probabilities_no_measurement.png}\\[3pt]
\includegraphics[width=.49\textwidth]{X_lock_line.png}\hfill
\includegraphics[width=.49\textwidth]{X_lock_circle.png}
\caption{Branch node: geometry; ensemble readout without measurement; channel probabilities; and readouts after dephasing lock to (left) line and (right) circle.}
\label{fig:branch}
\end{figure}

\section{Superdeterminism as Temporal Synchronization}
\subsection{Hidden variable and deterministic outcomes}
Let the hidden variable be a global \emph{clock phase} $\lambda\in[0,2\pi)$. Outcomes for settings $a,b$ are local and deterministic:
\begin{equation}
A(a,\lambda)=\mathrm{sgn}\!\big(\cos(\lambda-a)\big),\qquad
B(b,\lambda)=-\mathrm{sgn}\!\big(\cos(\lambda-b)\big).
\end{equation}

\subsection{Measurement independence is violated}
The setting generators are themselves functions of synchronized envelope structure; therefore $\rho(\lambda\,|\,a,b)\neq \rho(\lambda)$. Choosing a piecewise-uniform $\rho(\lambda\,|\,a,b)$ concentrated on sectors where $AB=\pm 1$ exactly reproduces the quantum correlator $E(a,b)=-\cos|a-b|$.

\subsection{Desynchronization knob and CHSH}
Let $\varepsilon\in[0,1]$ be the fraction of trials where $\lambda$ is drawn from an \emph{independent} uniform distribution (no sync). Then
\begin{equation}
S(\varepsilon)=(1-\varepsilon)\,2\sqrt2+2\,\varepsilon.
\end{equation}
Monte Carlo with $N=2\!\times\!10^5$ samples per pair reproduces this trend (fig.~\ref{fig:chsh}).

\begin{figure}[h!]
\centering
\includegraphics[width=.7\textwidth]{chsh_S_vs_epsilon.png}\\[6pt]
\includegraphics[width=.7\textwidth]{chsh_S_vs_delta_t.png}
\caption{Top: CHSH $S$ vs synchronization loss $\varepsilon$. Bottom: $S$ vs desynchronization time $\Delta t$ for $\varepsilon(\Delta t)=1-e^{-\Delta t/\tau}$ (here $\tau=5$).}
\label{fig:chsh}
\end{figure}

\section{Information Bounds and No-Signalling}
Even with global envelope correlations, accessible mutual information across a partition $A|B$ is bottlenecked by the envelope-induced sparsity of $\Keff$:
\begin{equation}
I_\text{accessible}(A\!:\!B)\;\le\; C_\text{cut}(A|B)=\sum_{e\in \mathrm{min\;cut}} C_e,
\end{equation}
where $C_e$ are link capacities induced by the filtered couplings. This bound underwrites no-signalling while allowing strong correlations.

\section{Retrocausality without Paradoxes}
Let $\mathcal E_{+}$ and $\mathcal E_{-}$ be forward/backward CP maps induced by envelopes. On the einselected subspace they commute and form a consistent two-time boundary problem:
\begin{equation}
[\mathcal E_{+},\mathcal E_{-}]=0,\qquad
\mathcal E_{-}\circ \mathcal E_{+}(\rho)=\rho.
\end{equation}
Retro-influences restrict histories but do not generate contradictions.

\section{Predictions and Experimental Knobs}
\begin{enumerate}[label=(\alph*)]
\item \textbf{Synchronization vs dephasing sweep:} Increasing synchronization (lower $\varepsilon$) and lowering dephasing increases $S$ toward $2\sqrt2$; the reverse trend lowers $S$ toward $2$.
\item \textbf{Indefinite causal order collapse:} As dephasing grows, causal witnesses in quantum-switch experiments should move monotonically into the causally separable region.
\item \textbf{Entanglement--geometry knob:} Varying envelope strength that sparsifies $\Keff$ should reduce operational distance and entanglement across a cut.
\end{enumerate}

\section{Conclusion}
TFE provides a small set of primitives that produce emergent space, constant intrinsic-speed propagation along timelines, measurement as phase locking, mass as gap/confinement/mixing, and a superdeterministic variant via temporal synchronization that preserves locality and no-signalling. The simulations illustrate concrete, falsifiable knobs. Future work includes a field-theoretic coarse graining, detailed capacity bounds for $\Keff$, and laboratory tests that co-vary synchronization and dephasing.

\section*{Acknowledgments}
We thank collaborators in spirit; any errors are ours.

\begin{thebibliography}{10}\setlength{\itemsep}{2pt}
\bibitem{Zurek2003} W. H. Zurek, ``Decoherence, einselection, and the quantum origins of the classical,'' \emph{Rev. Mod. Phys.} \textbf{75}, 715 (2003).
\bibitem{Maldacena1999} J. Maldacena, ``The large-$N$ limit of superconformal field theories and supergravity,'' \emph{Int. J. Theor. Phys.} \textbf{38}, 1113 (1999).
\bibitem{RyuTakayanagi2006} S. Ryu and T. Takayanagi, ``Holographic derivation of entanglement entropy,'' \emph{Phys. Rev. Lett.} \textbf{96}, 181602 (2006).
\bibitem{MaldacenaSusskind2013} J. Maldacena and L. Susskind, ``Cool horizons for entangled black holes,'' \emph{Fortschr. Phys.} \textbf{61}, 781 (2013).
\bibitem{ABL1964} Y. Aharonov, P. G. Bergmann, and J. L. Lebowitz, ``Time Symmetry in the Quantum Process,'' \emph{Phys. Rev.} \textbf{134}, B1410 (1964).
\bibitem{Oreshkov2012} O. Oreshkov, F. Costa, and \v{C}. Brukner, ``Quantum correlations with no causal order,'' \emph{Nat. Commun.} \textbf{3}, 1092 (2012).
\bibitem{LiebRobinson1972} E. H. Lieb and D. W. Robinson, ``The finite group velocity of quantum spin systems,'' \emph{Commun. Math. Phys.} \textbf{28}, 251 (1972).
\bibitem{Pastawski2015} F. Pastawski, B. Yoshida, D. Harlow, and J. Preskill, ``Holographic quantum error-correcting codes,'' \emph{JHEP} \textbf{06}, 149 (2015).
\bibitem{SachdevYe1993} S. Sachdev and J. Ye, ``Gapless spin-fluid ground state in a random, quantum Heisenberg magnet,'' \emph{Phys. Rev. Lett.} \textbf{70}, 3339 (1993).
\bibitem{MaldacenaStanford2016} J. Maldacena and D. Stanford, ``Remarks on the Sachdev-Ye-Kitaev model,'' \emph{Phys. Rev. D} \textbf{94}, 106002 (2016).
\end{thebibliography}

\end{document}
